% Section 2. Drafted from docs/PAPER-OUTLINE.md section 2 and the Coq definitions (GovernanceConverse.v,
% CausalReplay.v, RhoStar.v).

\section{Setting}\label{sec:setting}

\paragraph{Registries.} A \emph{registry} fixes a set of states $X$, a validity predicate $\valid$ (the
business rules), a set of events, and a repair. Each event $e$ has a raw transition $e : X \to X$, which may
leave the valid states. The repair is given by a step $\rho : X \to X$ applied to invalid states, with a
potential $\Phi : X \to \mathbb{N}$ that strictly decreases under $\rho$ on invalid states, so that repair
terminates. This termination condition is called WFC. Iterating $\rho$ until the state is valid gives the
repair map $N = \rho^{*}$, and the governed step of an event is $\gov(e) = \rho^{*} \circ e$.

\paragraph{Executions.} A replica holds a configuration: a state and a buffer of pending events. A step
either applies a pending event or, if the state is invalid, applies $\rho$. A normal form is a valid state
with an empty buffer. A state is reachable from a start $s_0$ when some sequence of event and repair steps
leads to it; the exact conditions of the paper quantify over reachable states, and the counterexamples of
\cref{sec:main-single} show that this restriction cannot be dropped.

\paragraph{Delivery models.} Which orders a replica may see is fixed by a delivery model: free delivery
(any order), causal delivery (orders consistent with a happens-before relation), at-least-once delivery
(also duplicates), and declared independence (only declared pairs may be swapped). \Cref{sec:layer-a}
treats them uniformly.

\paragraph{Networks of registries.} A federation connects registries by transport maps along which values
propagate; a network may be acyclic or contain cycles, and propagation may be lossy. These are the
settings of layer~B and of most of the regime map; their definitions are introduced where they are used
(\cref{sec:layer-b,sec:map}).

\paragraph{Mechanization.} Every result named in the paper is a theorem in a Coq development, cited by its
name in typewriter font. Every theorem is closed under the global context, with no axioms and no
admitted proofs, and continuous integration checks this on three versions of Coq on every change. The
statements in the paper are read from the source; where the paper's notation differs from the source,
the content does not.

\begin{center}\small
\begin{tabular}{@{}ll@{}}
\toprule
Notation & Meaning \\
\midrule
$X$, $\valid$ & states and the validity predicate \\
$e : X \to X$ & the raw transition of an event \\
$\rho$, $\Phi$ & a repair step and its potential (WFC: $\Phi$ decreases on invalid states) \\
$N = \rho^{*}$ & the repair map: apply $\rho$ until valid \\
$\gov(e) = N \circ e$ & the governed step \\
$\reach(s_0, \sigma)$ & $\sigma$ is reachable from $s_0$ \\
CC1, CC2 & the two commutation conditions of \cref{thm:cc} \\
\bottomrule
\end{tabular}
\end{center}
